\documentclass{article}
\usepackage{amsmath,amssymb,amsthm}
\usepackage{algorithm,algpseudocode}
\usepackage{geometry}
\geometry{margin=1in}
\newtheorem{theorem}{Theorem}
\newtheorem{lemma}{Lemma}
\newtheorem{assumption}{Assumption}
\newtheorem{corollary}{Corollary}
\newcommand{\prox}{\operatorname{prox}}
\newcommand{\R}{\mathbb{R}}

\begin{document}

\section*{Proximal Gradient Method (Forward–Backward Splitting)}

\section*{Mathematical Formulation}
Let 
\[
F(x)=f(x)+g(x),\qquad x\in\R^{d},
\]
where  

\begin{itemize}
    \item $f:\R^{d}\!\to\!\R$ is convex and differentiable with $L$‑Lipschitz continuous gradient:
    \[
    \|\nabla f(u)-\nabla f(v)\|\le L\|u-v\|,\qquad\forall u,v\in\R^{d};
    \]
    \item $g:\R^{d}\!\to\!\R\cup\{+\infty\}$ is proper, closed and convex (possibly nonsmooth).
\end{itemize}
For a stepsize $\eta>0$, the proximal operator of $g$ is
\[
\prox_{\eta g}(z):=\arg\min_{x\in\R^{d}}\Big\{ g(x)+\frac{1}{2\eta}\|x-z\|^{2}\Big\}.
\]

The algorithm generates the sequence $\{x_{k}\}_{k\ge0}$ by
\begin{equation}
\boxed{\,x_{k+1}= \prox_{\eta g}\!\big(x_{k}-\eta\nabla f(x_{k})\big)\,}\tag{1}
\end{equation}
with an initial point $x_{0}\in\R^{d}$ and a stepsize $\eta\in(0,1/L]$.

\section*{Pseudocode}
\begin{algorithm}[H]
\caption{Proximal Gradient Method}
\begin{algorithmic}[1]
\Require Initial point $x_{0}$, stepsize $\eta\in(0,1/L]$, maximum iterations $K$.
\For{$k=0,1,\dots,K-1$}
    \State Compute gradient $g_{k}\gets \nabla f(x_{k})$.
    \State Take a forward step $y_{k}\gets x_{k}-\eta g_{k}$.
    \State Apply the proximal map $x_{k+1}\gets \prox_{\eta g}(y_{k})$.
\EndFor
\State \textbf{return} $x_{K}$.
\end{algorithmic}
\end{algorithm}

\section*{Dry Run Example}
Consider the scalar problem 
\[
f(w)=(w-3)^{2},\qquad g\equiv0,
\]
so that $\nabla f(w)=2(w-3)$ and $\prox_{\eta g}= \operatorname{id}$.  
Take $w_{0}=0$, stepsize $\eta=0.1$, and run three iterations.

\begin{center}
\begin{tabular}{c|c|c|c}
$k$ & $w_{k}$ & $\nabla f(w_{k})$ & $w_{k+1}=w_{k}-\eta\nabla f(w_{k})$\\ \hline
0 & $0$ & $-6$ & $0-0.1(-6)=0.6$\\
1 & $0.6$ & $-4.8$ & $0.6-0.1(-4.8)=1.08$\\
2 & $1.08$ & $-3.84$ & $1.08-0.1(-3.84)=1.464$\\
\end{tabular}
\end{center}
The objective values $F(w_{k})=f(w_{k})$ are
\[
f(0)=9,\; f(0.6)=5.76,\; f(1.08)=3.6864,\; f(1.464)=2.2817,
\]
demonstrating monotone decrease toward the optimum $w^{*}=3$.

\newpage

\section*{Assumptions}
\begin{assumption}[Smoothness of $f$]\label{as:lip}
$f$ is convex and differentiable with $L$‑Lipschitz gradient:
\[
\|\nabla f(u)-\nabla f(v)\|\le L\|u-v\|,\quad\forall u,v\in\R^{d}.
\]
\end{assumption}
\begin{assumption}[Convexity of $g$]\label{as:convex}
$g$ is proper, closed and convex (possibly nonsmooth).
\end{assumption}
\begin{assumption}[Stepsize]\label{as:stepsize}
The constant stepsize satisfies $0<\eta\le 1/L$.
\end{assumption}
\begin{assumption}[Existence of a minimiser]\label{as:opt}
The composite problem $\min_{x\in\R^{d}}F(x)$ admits at least one solution $x^{\star}$.
\end{assumption}

\section*{Main Convergence Theorem}
\begin{theorem}[Ergodic $O(1/k)$ rate]\label{thm:rate}
Under Assumptions \ref{as:lip}–\ref{as:opt}, the iterates generated by (1) satisfy for every $k\ge0$
\[
F(x_{k})-F(x^{\star})\le \frac{\|x_{0}-x^{\star}\|^{2}}{2\eta\,k}.
\]
Consequently,
\[
\min_{0\le i\le k-1}\bigl(F(x_{i})-F(x^{\star})\bigr)=O\!\left(\frac{1}{k}\right).
\]
\end{theorem}

\section*{Theoretical Properties}
\begin{itemize}
    \item \textbf{Stability.} The non‑expansiveness of $\prox_{\eta g}$ ensures that the algorithm never amplifies errors: 
    $\|x_{k+1}-x^{\star}\|\le\|x_{k}-\eta\nabla f(x_{k})- (x^{\star}-\eta\nabla f(x^{\star}))\|$.
    \item \textbf{Stepsize sensitivity.} Any $\eta\in(0,1/L]$ guarantees the descent inequality used in the proof; larger $\eta$ may break the contraction.
    \item \textbf{Comparison.} When $g\equiv0$, (1) reduces to plain gradient descent and Theorem~\ref{thm:rate} coincides with the classical $O(1/k)$ bound. For nonsmooth $g$, the same rate holds without any smoothing of $g$.
\end{itemize}

\section*{Discussion}
The theorem provides a deterministic guarantee; stochastic variants require additional variance bounds. The $O(1/k)$ rate is optimal for first‑order methods on the class of convex composite problems with Lipschitz smooth part.

\newpage

\section*{Preliminaries (Definitions and Lemmas)}

\begin{lemma}[Descent Lemma]\label{lem:descent}
If $\nabla f$ is $L$‑Lipschitz, then for any $u,v\in\R^{d}$,
\[
f(v)\le f(u)+\langle \nabla f(u),v-u\rangle+\frac{L}{2}\|v-u\|^{2}.
\]
\end{lemma}
\begin{proof}
Standard consequence of the integral form of the mean value theorem and the Lipschitz property.
\end{proof}

\begin{lemma}[Firm non‑expansiveness of the proximal map]\label{lem:prox}
For any $u,v\in\R^{d}$ and any $\eta>0$,
\[
\|\prox_{\eta g}(u)-\prox_{\eta g}(v)\|^{2}
\le \langle u-v,\prox_{\eta g}(u)-\prox_{\eta g}(v)\rangle .
\]
\end{lemma}
\begin{proof}
The optimality condition of the proximal subproblem yields
\[
0\in\partial g(\prox_{\eta g}(u))+\frac{1}{\eta}\bigl(\prox_{\eta g}(u)-u\bigr), 
\qquad
0\in\partial g(\prox_{\eta g}(v))+\frac{1}{\eta}\bigl(\prox_{\eta g}(v)-v\bigr).
\]
Subtracting, using monotonicity of $\partial g$, and rearranging gives the stated inequality.
\end{proof}

\section*{Main Proof (Step‑by‑Step Derivation)}

Let $x^{\star}$ be any minimiser of $F$. Define the auxiliary point
\[
y_{k}=x_{k}-\eta\nabla f(x_{k}) .
\]
By the algorithm, $x_{k+1}= \prox_{\eta g}(y_{k})$.

\bigskip
\noindent\textbf{[Step 1]}  \textit{Optimality condition of the proximal step}.  
From Lemma~\ref{lem:prox} with $u=y_{k}$ and $v=x^{\star}$,
\[
\|x_{k+1}-x^{\star}\|^{2}
\le \langle y_{k}-x^{\star},\,x_{k+1}-x^{\star}\rangle .
\tag{2}
\]

\noindent\textbf{[Step 2]}  \textit{Insert the definition of $y_{k}$}.  
Since $y_{k}=x_{k}-\eta\nabla f(x_{k})$,
\[
\langle y_{k}-x^{\star},\,x_{k+1}-x^{\star}\rangle
= \langle x_{k}-x^{\star},\,x_{k+1}-x^{\star}\rangle
-\eta\langle \nabla f(x_{k}),\,x_{k+1}-x^{\star}\rangle .
\tag{3}
\]

\noindent\textbf{[Step 3]}  \textit{Rewrite the first inner product}.  
Using the identity $2\langle a,b\rangle =\|a\|^{2}+\|b\|^{2}-\|a-b\|^{2}$ with $a=x_{k}-x^{\star}$ and $b=x_{k+1}-x^{\star}$,
\[
2\langle x_{k}-x^{\star},\,x_{k+1}-x^{\star}\rangle
= \|x_{k}-x^{\star}\|^{2}+\|x_{k+1}-x^{\star}\|^{2}
-\|x_{k+1}-x_{k}\|^{2}.
\tag{4}
\]

\noindent\textbf{[Step 4]}  \textit{Combine (2)–(4)}.  
Multiplying (2) by $2$ and substituting (3) and (4) yields
\[
\begin{aligned}
2\|x_{k+1}-x^{\star}\|^{2}
&\le \|x_{k}-x^{\star}\|^{2}+\|x_{k+1}-x^{\star}\|^{2}
-\|x_{k+1}-x_{k}\|^{2}\\
&\quad -2\eta\langle \nabla f(x_{k}),\,x_{k+1}-x^{\star}\rangle .
\end{aligned}
\]
Rearranging,
\[
\|x_{k+1}-x^{\star}\|^{2}
\le \|x_{k}-x^{\star}\|^{2}
-\|x_{k+1}-x_{k}\|^{2}
-2\eta\langle \nabla f(x_{k}),\,x_{k+1}-x^{\star}\rangle .
\tag{5}
\]

\noindent\textbf{[Step 5]}  \textit{Bound the smooth part via the descent lemma}.  
Applying Lemma~\ref{lem:descent} with $u=x_{k}$ and $v=x_{k+1}$,
\[
f(x_{k+1})\le f(x_{k})+\langle \nabla f(x_{k}),\,x_{k+1}-x_{k}\rangle
+\frac{L}{2}\|x_{k+1}-x_{k}\|^{2}.
\tag{6}
\]

\noindent\textbf{[Step 6]}  \textit{Add $g(x_{k+1})$ to both sides of (6)} and use the definition $F=f+g$:
\[
F(x_{k+1})\le f(x_{k})+g(x_{k+1})+\langle \nabla f(x_{k}),\,x_{k+1}-x_{k}\rangle
+\frac{L}{2}\|x_{k+1}-x_{k}\|^{2}.
\tag{7}
\]

\noindent\textbf{[Step 7]}  \textit{Use the proximal optimality inequality}.  
From the definition of $\prox_{\eta g}$,
\[
g(x_{k+1})+\frac{1}{2\eta}\|x_{k+1}-y_{k}\|^{2}
\le g(x^{\star})+\frac{1}{2\eta}\|x^{\star}-y_{k}\|^{2}.
\tag{8}
\]
Recall $y_{k}=x_{k}-\eta\nabla f(x_{k})$, hence
\[
\|x_{k+1}-y_{k}\|^{2}= \|x_{k+1}-x_{k}+\eta\nabla f(x_{k})\|^{2}
= \|x_{k+1}-x_{k}\|^{2}+2\eta\langle \nabla f(x_{k}),\,x_{k+1}-x_{k}\rangle+\eta^{2}\|\nabla f(x_{k})\|^{2},
\]
and similarly for the term involving $x^{\star}$. Substituting these expansions into (8) and simplifying gives
\[
g(x_{k+1})\le g(x^{\star})+\langle \nabla f(x_{k}),\,x^{\star}-x_{k+1}\rangle
+\frac{1}{2\eta}\bigl(\|x^{\star}-x_{k}\|^{2}-\|x^{\star}-x_{k+1}\|^{2}
-\|x_{k+1}-x_{k}\|^{2}\bigr).
\tag{9}
\]

\noindent\textbf{[Step 8]}  \textit{Combine (7) and (9)}.  Adding (9) to (7) eliminates the term $\langle \nabla f(x_{k}),x_{k+1}-x_{k}\rangle$:
\[
\begin{aligned}
F(x_{k+1}) &\le f(x_{k})+g(x^{\star})
+\langle \nabla f(x_{k}),\,x^{\star}-x_{k}\rangle
+\frac{1}{2\eta}\bigl(\|x^{\star}-x_{k}\|^{2}
-\|x^{\star}-x_{k+1}\|^{2}
-\|x_{k+1}-x_{k}\|^{2}\bigr) \\
&\quad +\frac{L}{2}\|x_{k+1}-x_{k}\|^{2}.
\end{aligned}
\tag{10}
\]

\noindent\textbf{[Step 9]}  \textit{Use convexity of $f$}.  
Convexity yields
\[
f(x_{k})+ \langle \nabla f(x_{k}),\,x^{\star}-x_{k}\rangle \le f(x^{\star}).
\tag{11}
\]

\noindent\textbf{[Step 10]}  \textit{Insert (11) into (10)} and collect the quadratic terms:
\[
F(x_{k+1})-F(x^{\star})
\le \frac{1}{2\eta}\bigl(\|x^{\star}-x_{k}\|^{2}
-\|x^{\star}-x_{k+1}\|^{2}\bigr)
+\Bigl(\frac{L}{2}-\frac{1}{2\eta}\Bigr)\|x_{k+1}-x_{k}\|^{2}.
\tag{12}
\]

\noindent\textbf{[Step 11]}  \textit{Choose stepsize $\eta\le 1/L$}.  
With $0<\eta\le 1/L$, the coefficient in front of $\|x_{k+1}-x_{k}\|^{2}$ is non‑positive:
\[
\frac{L}{2}-\frac{1}{2\eta}\le 0.
\]
Hence we can drop the last term, obtaining the fundamental descent inequality
\[
F(x_{k+1})-F(x^{\star})
\le \frac{1}{2\eta}\bigl(\|x^{\star}-x_{k}\|^{2}
-\|x^{\star}-x_{k+1}\|^{2}\bigr).
\tag{13}
\]

\noindent\textbf{[Step 12]}  \textit{Telescoping sum}.  
Sum (13) for $k=0,\dots, K-1$:
\[
\sum_{k=0}^{K-1}\bigl(F(x_{k+1})-F(x^{\star})\bigr)
\le \frac{1}{2\eta}\bigl(\|x^{\star}-x_{0}\|^{2}
-\|x^{\star}-x_{K}\|^{2}\bigr)
\le \frac{\|x_{0}-x^{\star}\|^{2}}{2\eta}.
\tag{14}
\]

\noindent\textbf{[Step 13]}  \textit{Derive the $O(1/K)$ bound}.  
Since the left‑hand side of (14) is at least $K\min_{0\le k\le K-1}\bigl(F(x_{k+1})-F(x^{\star})\bigr)$,
\[
\min_{0\le k\le K-1}\bigl(F(x_{k+1})-F(x^{\star})\bigr)
\le \frac{\|x_{0}-x^{\star}\|^{2}}{2\eta K}.
\tag{15}
\]
This proves Theorem~\ref{thm:rate}. $\square$

\section*{Rate Derivation (With Constants)}
From (13) we have for each $k$,
\[
F(x_{k})-F(x^{\star})\le \frac{\|x_{0}-x^{\star}\|^{2}}{2\eta\,k},
\]
where the constant $C:=\|x_{0}-x^{\star}\|^{2}/(2\eta)$ depends only on the initial distance and the stepsize. The bound is tight for the worst‑case convex quadratic $f$ with $g\equiv0$.

\section*{Special Cases}
\begin{itemize}
    \item \textbf{Pure gradient descent.} If $g\equiv0$, then $\prox_{\eta g}=\operatorname{id}$ and (13) reduces to the classical gradient‑descent inequality.
    \item \textbf{Indicator function.} If $g=\iota_{\mathcal{C}}$ for a closed convex set $\mathcal{C}$, then $\prox_{\eta g}$ becomes the Euclidean projection onto $\mathcal{C}$, and the method coincides with the projected gradient method; the same $O(1/k)$ rate holds.
    \item \textbf{Strongly convex case.} If $F$ is $\mu$‑strongly convex, a refined analysis (omitted for brevity) yields a linear rate $F(x_{k})-F(x^{\star})\le (1-\eta\mu)^{k}\frac{\|x_{0}-x^{\star}\|^{2}}{2\eta}$.
\end{itemize}

\end{document}